% Source problem: https://www.ap.t.u-tokyo.ac.jp/themes/ap_t_u-tokyo/graduate/pdf/h28ap_exam.pdf
\documentclass[11pt,a4paper]{article}
\usepackage[margin=24mm]{geometry}
\usepackage{fontspec}
\usepackage{xeCJK}
\setCJKmainfont{HaranoAjiMincho-Regular.otf}
\setCJKsansfont{HaranoAjiGothic-Medium.otf}
\usepackage{amsmath,amssymb}
\usepackage{tikz}
\usetikzlibrary{arrows.meta}
\setlength{\parindent}{1em}
\setlength{\parskip}{0.35em}
\renewcommand{\figurename}{図}
\title{物理工学専攻 院試解答集\\2016年度 物理学II 第3問}
\author{政岡凜太郎}
\date{}
\begin{document}
\maketitle

\section*{[1] 束縛電子の運動}
電子の電荷は$-e$なので、電場による力は$-e\boldsymbol E$である。共振振動数$\omega_0$の調和復元力も加えると、
\begin{equation}
  \boxed{m\ddot{\boldsymbol x}=-m\omega_0^2\boldsymbol x-e\boldsymbol E.}
\end{equation}

\section*{[2] 分極の運動方程式}
数密度$N$の電子による分極は$\boldsymbol P=-eN\boldsymbol x$である。これを時間で2回微分し、[1]を使うと
\begin{equation}
  \boxed{\ddot{\boldsymbol P}+\omega_0^2\boldsymbol P
  =\frac{Ne^2}{m}\boldsymbol E.}
\end{equation}
右辺の符号が正なのは、$\boldsymbol P$と電子変位$\boldsymbol x$が反対向きだからである。

\section*{[3] 誘電率}
電場と分極が角周波数$\omega$で振動するので、$\ddot P=-\omega^2P$である。従って
\begin{equation}
  (\omega_0^2-\omega^2)P=\frac{Ne^2}{m}E,
  \qquad
  \boxed{P=\frac{Ne^2}{m(\omega_0^2-\omega^2)}E}.
\end{equation}
マクロな関係$D=\varepsilon_0E+P=\varepsilon(\omega)E$より
\begin{equation}
  \varepsilon(\omega)=\varepsilon_0+
  \frac{Ne^2}{m(\omega_0^2-\omega^2)}.
\end{equation}

\section*{[4] 誘電体中の分散関係}
問題文の波動方程式(3)に平面波を代入すると$-k^2E+\varepsilon(\omega)\mu_0\omega^2E=0$だから
\begin{equation}
  \boxed{k^2=\mu_0\varepsilon(\omega)\omega^2}.
\end{equation}
プラズマ角周波数を$\omega_p^2=Ne^2/(\varepsilon_0m)$、
$\omega_1=\sqrt{\omega_0^2+\omega_p^2}$とおく。真空中の光速$c=1/\sqrt{\varepsilon_0\mu_0}$を使えば
\begin{equation}
  \boxed{k^2=\frac{\omega^2}{c^2}
  \frac{\omega_1^2-\omega^2}{\omega_0^2-\omega^2}}.
\end{equation}

\section*{[5] 誘電体の実部・虚部}
正の角周波数に対して、分子と分母の符号を比べる。$+z$方向へ進む波には、伝搬域で$\operatorname{Re}k>0$、非伝搬域で$\operatorname{Im}k>0$となる枝を選ぶ。
\begin{equation}
  \begin{array}{c|c|c}
    \text{周波数域}&\operatorname{Re}k&\operatorname{Im}k\\ \hline
    0<\omega<\omega_0&
    \displaystyle\frac{\omega}{c}\sqrt{\frac{\omega_1^2-\omega^2}{\omega_0^2-\omega^2}}&0\\[1.1em]
    \omega_0<\omega<\omega_1&0&
    \displaystyle\frac{\omega}{c}\sqrt{\frac{\omega_1^2-\omega^2}{\omega^2-\omega_0^2}}\\[1.1em]
    \omega>\omega_1&
    \displaystyle\frac{\omega}{c}\sqrt{\frac{\omega^2-\omega_1^2}{\omega^2-\omega_0^2}}&0
  \end{array}
\end{equation}
誘電体の非伝搬帯では$k=i\kappa$（$\kappa>0$）を選ぶ。位相因子は$e^{ikz}=e^{-\kappa z}$となり、$+z$方向へ進むと振幅が減衰する。これはエバネッセント場の空間的減衰であり、この無減衰模型での吸収損失を意味しない。共振$\omega_0$では理想化した無減衰模型の$k$が発散し、$\omega_1$では$k=0$となる。

図の例では $\Omega=\omega/\omega_0$、$K=ck/\omega_0$ と規格化し、$\omega_1/\omega_0=2$ を選ぶ。従って $\omega_p/\omega_0=\sqrt{3}$ であり、実部・虚部の枝は直前の分散式から
\begin{equation}
 \operatorname{Re}K=\begin{cases}
 \displaystyle\Omega\sqrt{\frac{4-\Omega^2}{1-\Omega^2}},&0<\Omega<1,\\
 0,&1<\Omega<2,\\
 \displaystyle\Omega\sqrt{\frac{\Omega^2-4}{\Omega^2-1}},&\Omega>2,
 \end{cases}\quad
 \operatorname{Im}K=\begin{cases}
 \displaystyle\Omega\sqrt{\frac{4-\Omega^2}{\Omega^2-1}},&1<\Omega<2,\\
 0,&0<\Omega<1\text{ または }\Omega>2.
 \end{cases}
\end{equation}

\begin{figure}[htbp]
\centering
\begin{tikzpicture}[>=Stealth,font=\small,x=1.25cm,y=0.8cm]
  \begin{scope}
    \draw[->] (0,0)--(3.25,0) node[right] {$\Omega$};
    \draw[->] (0,0)--(0,1.9) node[above] {$\operatorname{Re}K$};
    \draw[dashed] (1,0)--(1,1.65);
    \draw[dashed] (2,0)--(2,1.55);
    \begin{scope}
    \clip (0,0) rectangle (3.1,1.8);
    \draw[thick] plot[domain=0:0.985,samples=100] (\x,{\x*sqrt((4-\x^2)/(1-\x^2))});
    \draw[thick] plot[domain=2:3.1,samples=100] (\x,{\x*sqrt((\x^2-4)/(\x^2-1))});
    \end{scope}
    \node[below] at (1,0) {$1$};
    \node[below] at (2,0) {$2$};
    \node at (1.5,0.3) {$0$};
  \end{scope}
  \begin{scope}[xshift=5.2cm]
    \draw[->] (0,0)--(3.15,0) node[right] {$\Omega$};
    \draw[->] (0,0)--(0,1.9) node[above] {$\operatorname{Im}K$};
    \draw[dashed] (1,0)--(1,1.65);
    \draw[dashed] (2,0)--(2,1.55);
    \begin{scope}
    \clip (0,0) rectangle (3.1,1.8);
    \draw[thick,dashed] plot[domain=1.015:2,samples=100] (\x,{\x*sqrt((4-\x^2)/(\x^2-1))});
    \end{scope}
    \node[below] at (1,0) {$1$};
    \node[below] at (2,0) {$2$};
  \end{scope}
\end{tikzpicture}
\caption{誘電体の規格化波数 $K$ の実部（実線）と虚部（破線）。共振近傍で発散する枝は枠外へ続く。}
\label{fig:dielectric-k}
\end{figure}

\section*{[6] 金属中の自由電子}
金属中の電子には束縛復元力がない。磁気力を無視すると、電場による運動方程式は
\begin{equation}
  \boxed{m\ddot{\boldsymbol x}=-e\boldsymbol E.}
\end{equation}
分極$\boldsymbol P=-en_e\boldsymbol x$を導入すれば、調和応答では
\begin{equation}
  P=-\frac{n_e e^2}{m\omega^2}E,\qquad
  \varepsilon_m(\omega)=\varepsilon_0\left(1-\frac{\omega_{pm}^2}{\omega^2}\right),
  \qquad \omega_{pm}^2=\frac{n_e e^2}{\varepsilon_0m}.
\end{equation}

\section*{[7] 金属の実部・虚部}
分散関係は
\begin{equation}
  k^2=\mu_0\varepsilon_0(\omega^2-\omega_{pm}^2)
  =\frac{\omega^2-\omega_{pm}^2}{c^2}.
\end{equation}
ここでは $\Omega=\omega/\omega_{pm}$、$K=ck/\omega_{pm}$ とおくと、正周波数で
\begin{equation}
 \operatorname{Re}K=\begin{cases}0,&0<\Omega<1,\\
 \sqrt{\Omega^2-1},&\Omega>1,\end{cases}\qquad
 \operatorname{Im}K=\begin{cases}\sqrt{1-\Omega^2},&0<\Omega<1,\\
 0,&\Omega>1.\end{cases}
\end{equation}
となる。
従って正の角周波数では
\begin{equation}
  \boxed{\begin{aligned}
    0<\omega<\omega_{pm}:&\quad \operatorname{Re}k=0,\quad
      \operatorname{Im}k=\frac1c\sqrt{\omega_{pm}^2-\omega^2},\\
    \omega>\omega_{pm}:&\quad \operatorname{Re}k=\frac1c\sqrt{\omega^2-\omega_{pm}^2},\quad
      \operatorname{Im}k=0.
  \end{aligned}}
\end{equation}
\begin{figure}[htbp]
\centering
\begin{tikzpicture}[>=Stealth,font=\small,x=1.25cm,y=0.8cm]
  \draw[->] (0,0)--(3.3,0) node[right] {$\Omega=\omega/\omega_{pm}$};
  \draw[->] (0,0)--(0,3.2) node[above] {$K=ck/\omega_{pm}$};
  \draw[dashed] (1,0)--(1,1.65);
  \draw[thick] plot[domain=1:3.1,samples=100] (\x,{sqrt(\x^2-1)});
  \draw[thick,dashed] plot[domain=0:1,samples=100] (\x,{sqrt(1-\x^2)});
  \node[below] at (1,0) {$1$};
  \node at (2.5,0.7) {$\operatorname{Re}K$};
  \node at (0.35,1.5) {$\operatorname{Im}K$};
\end{tikzpicture}
\caption{金属の規格化波数 $K$ の実部（実線）と虚部（破線）。}
\label{fig:metal-k}
\end{figure}

\section*{[8] 誘電体と金属の光学的性質}
可視光を各模型の低周波数側とみなす。束縛電子をもつ誘電体で$\omega\ll\omega_0$なら$k$は実数であり、波は媒質内を伝搬する。対して金属で$\omega<\omega_{pm}$なら$k$は虚数であり、電場は表面から内部へ指数関数的に減衰する。内部にはエバネッセント場が存在する。ここで$H_y$は磁場の強さの$y$成分である。位相因子$e^{i(kz-\omega t)}$に対するFaradayの法則$\nabla\times\boldsymbol E=-\mu_0\partial_t\boldsymbol H$の$y$成分は$ikE_x=i\omega\mu_0H_y$だから、$H_y=kE_x/(\mu_0\omega)$となる。$k$が純虚数なので、時間平均の法線エネルギー流は$\operatorname{Re}(E_xH_y^*)/2=|E_x|^2\operatorname{Re}(k)/(2\mu_0\omega)=0$である。散逸のない半無限金属に平均電力は伝わらず、入射電力は反射される。金属内部の空間的減衰を吸収損失とは解釈しない。

\end{document}
